%\nonstopmode
%\errorstopmode  %% opposite of nonstopmode
\documentclass[t,fleqn,usenames,dvipsnames]{beamer}
%\documentclass{article}\usepackage{beamerarticle}

\makeatletter
%\setlength\@mathmargin{1em}  %% reduce math margin
%\def\leqn{\tagsleft@true}
%\def\reqn{\tagsleft@false}
\def\fleq{\@fleqntrue\let\mathindent\@mathmargin \@mathmargin=\leftmargini}
\def\cneq{\@fleqnfalse}
\g@addto@macro{\endsubequations}{\addtocounter{equation}{-1}}
\makeatother

\mode<presentation>
{
  % \usetheme{AnnArbor}
  \useinnertheme[shadow=true]{rounded}
  \useoutertheme{infolines}
  %\setbeamercolor*{title in head/foot}{parent=palette secondary}
  %\useoutertheme{shadow}
  \usecolortheme{wolverine}

  \setbeamerfont{block title}{size={}}
  \setbeamercolor{titlelike}{parent=structure,bg=yellow!85!orange}

  % My modification: centered title with white background
  \setbeamercolor{frametitle}{bg=white}
  \setbeamertemplate{frametitle}{
      \Large\insertframetitle
      \par
      %{\centering\large\insertframesubtitle}
      %\par
  }
  % \setbeamertemplate{frametitle}{
  %   %\begin{center}
  %     \Large\insertframetitle
  %     \vspace{1ex}
  %     %\hline
  %     \par
  %   %\end{center}
  % }
  % %% Does not work
  % \setbeamercolor{frametitlecentered}{bg=white}
  % \setbeamertemplate{frametitlecentered}{
  %   \begin{center}
  %     \Large\insertframetitle
  %     \par
  %   \end{center}
  % }
  % plus red emphasis
  \setbeamercolor{alerted text}{fg=red!70!black}

  % More space between items
  %\defbeamertemplate*{itemize/enumerate body begin}{default}{\itemsep1ex}


  \setbeamercolor{math text}{parent=titlelike}
%  \setbeamercolor{math text displayed}{parent=palette primary}

  % Fix background of theorems/proof
  \setbeamercolor{block body}{parent=palette primary}%{bg=yellow!85!orange}
  \setbeamercolor{block title}{parent=palette secondary}%{bg=orange}



  \setbeamercovered{transparent}
  % or whatever (possibly just delete it)

} % end mode presentation

%\usepackage{tcolorbox}
\usepackage[english]{babel}
\usepackage[utf8x]{inputenc}
%\usepackage{times}  % DYSFUNCTIONAL! looks awful when mixing mathtt and ordinary math.
\usepackage{ifthen}
%\usepackage[T1]{fontenc} %no effect
% Or whatever. Note that the encoding and the font should match. If T1
% does not look nice, try deleting the line with the fontenc.
\usepackage{array} % for \newcolumntype
\usepackage{bbm}
\usepackage{calc} % for \widthof
\usepackage{amsmath}
\usepackage{amssymb}
\usepackage{stmaryrd}
\usepackage{alltt}
\usepackage[normalem]{ulem} % strikethrough \sout
\usepackage{cancel} % strikethrough math mode \cancel
%\usepackage{enumitem} % set label in itemize
\usepackage{pifont} % for \tickNo
%\usepackage{}
%\usepackage{listings}
\usepackage[all]{xy}
%\usepackage{proof}
%\usepackage{eurosym}
%\usepackage{graphics}
\usepackage{bibentry}
%\nobibliography{short}
\bibliographystyle{plain}
%\input{prooftree}

\newtheorem{takehomemessage}[theorem]{\textbf{Take home message}}

%\usepackage{agda}
%\AgdaNoSpaceAroundCode{}

\DeclareMathSymbol{\chkmark}{\mathord}{AMSa}{"58}

% RGB colors
\definecolor{darkred}{rgb}{0.5,0,0}
\definecolor{darkgreen}{rgb}{0,0.5,0}
\definecolor{darkblue}{rgb}{0,0,0.5}
\definecolor{dirtyred}{rgb}{0.7,0.2,0.1}
\definecolor{dirtygreen}{rgb}{0.2,0.4,0.1}
\definecolor{darkdirtygreen}{rgb}{0.13,0.25,0.07}
\definecolor{dirtyblue}{rgb}{0.07,0.2,0.5}
\definecolor{darkdirtyblue}{rgb}{0.1,0.15,0.35}
\definecolor{lightblue}{rgb}{0.5,0.5,1}
\definecolor{olivegreen}{rgb}{0.5,0.5,0}
\definecolor{brown}{rgb}{0.65,0.35,0} % almost gold
\definecolor{grey}{rgb}{0.33,0.33,0.33}
\definecolor{darkbrown}{rgb}{0.35,0.15,0}
\definecolor{darkgrey}{rgb}{0.16,0.16,0.16}

\newcommand{\textred}[1]{{\color{dirtyred}\textbf{#1}}}

%\newcommand{\tickYes}{{\color{darkgreen}\checkmark}} % does not work
  %since checkmark invokes math mode
\newcommand{\tickYes}{\ensuremath{\color{darkgreen}\chkmark}}
\newcommand{\tickNo}{{\color{dirtyred}\hspace{1pt}\ding{55}}}

%\usepackage[curve,matrix,arrow]{xy}
%\usepackage{tikz} % Drawing diagrams
%\usepackage{pgflibraryshapes} %Ellipses

\newcommand{\SKIP}[1]{}
\newcommand{\LESS}[1]{}

% MACROS:
\usepackage{agda}
\newcommand{\inst}{}
\input{colors}
\input{talkMacros}

% Semantics
\renewcommand{\den}[1]{\semcol{\lv} #1 \semcol{\rv}}
\newcommand{\posden}[1]{\poscol{\lv} #1 \poscol{\rv^{+}}}
\newcommand{\negden}[1]{\negcol{\lv} #1 \negcol{\rv^{-}}}
\newcommand{\cden}[1]{\den{\cxtcol{#1}}}
\renewcommand{\A}[1][]{\semcol{\mathcal{A}}_{\cxtcol{#1}}}
\renewcommand{\B}[1][]{\semcol{\mathcal{B}}_{\cxtcol{#1}}}
\renewcommand{\C}[1][]{\semcol{\mathcal{C}}_{\cxtcol{#1}}}
\renewcommand{\N}[1][]{\negcol{\mathcal{N}}_{\cxtcol{#1}}}
\newcommand{\Np}[1][]{\negcol{\mathcal{N'}}_{\cxtcol{#1}}}
\renewcommand{\P}[1][]{\poscol{\mathcal{P}}_{\cxtcol{#1}}}
\newcommand{\Pp}[1][]{\poscol{\mathcal{P'}}_{\cxtcol{#1}}}

\newcommand{\inject}[1]{\semcol{\iota_{#1}}}

\newcommand{\semrun}{\covcol{\mathsf{run}}}
\newcommand{\semcase}{\covcol{\mathsf{case}}}
\newcommand{\covcase}{\covcol{\mathsf{case}}}
\newcommand{\covabort}{\covcol{\mathsf{abort}}}
\newcommand{\covreturn}{\covcol{\mathsf{return}}}
\newcommand{\covjoin}{\covcol{\mathsf{join}}}
\newcommand{\covmon}{\covcol{\mathsf{mon}}^{\cover{}}}
\newcommand{\covmap}{\covcol{\mathsf{map}}}
\newcommand{\covsmap}{\covcol{\widehat{\mathsf{map}}}}

% \renewcommand{\tmc}[1]{\mathsf{#1}}
% \renewcommand{\tmcol}[1]{\mathsf{#1}}
\renewcommand{\varcol}{\tmcol}

\renewcommand{\Var}[2]{\tVar\;#1\;#2}
\renewcommand{\ru}[2]{\dfrac{%
  \begin{array}[b]{@{}l@{}} #1 \end{array}\hfill%
  }{%
  \begin{array}[t]{@{}l@{}} #2 \end{array}\hfill%
  }}

% black text in mbox
\newcommand{\mybox}[1]{\mbox{\color{black}#1}}

% meta-level logic
\renewcommand{\mfor}{\ \mybox{for}\ }
\renewcommand{\mforsome}{\ \mybox{for some}\ }
\renewcommand{\mthen}{\ \mybox{then}\ }
\renewcommand{\mif}{\ \mybox{if}\ }
\renewcommand{\miff}{\ \mybox{iff}\ }
\renewcommand{\motherwise}{\ \mybox{otherwise}}
\renewcommand{\mundefined}{\mybox{undefined}}
\renewcommand{\mnot}{\mybox{not}\ }
\renewcommand{\mand}{\ \mybox{and}\ }
\renewcommand{\mor}{\ \mybox{or}\ }
\renewcommand{\mimplies}{\ \mybox{implies}\ }
\renewcommand{\mimply}{\ \mybox{imply}\ }
\renewcommand{\mforall}{\ \mybox{for all}\ }
\renewcommand{\mexists}{\mybox{exists}\ }
\renewcommand{\mexist}{\mybox{exist}\ }
\renewcommand{\mtrue}{\mybox{true}}
\renewcommand{\mwith}{\ \mybox{with}\ }
\renewcommand{\mwhere}{\ \mybox{where}\ }
\renewcommand{\mholds}{\ \mybox{holds}\ }
\renewcommand{\munless}{\ \mybox{unless}\ }
\renewcommand{\mboth}{\ \mybox{both}\ }
\renewcommand{\msuchthat}{\ \mybox{such that}\ }
% proofs
\renewcommand{\msince}{\mybox{since}\ }
\renewcommand{\mdef}{\mybox{by def.}}
\renewcommand{\mass}{\mybox{assumption}}
\renewcommand{\mhyp}{\mybox{by hyp.}}
\renewcommand{\mlemma}[1]{\mybox{by Lemma~#1}}
\renewcommand{\mih}[1][]{\mybox{by ind.hyp.}#1}
\renewcommand{\mgoal}[1][]{\mybox{goal\ifthenelse{\equal{#1}{}}{}{~#1}}}
\renewcommand{\mby}[1]{\mybox{by #1}}
\renewcommand{\minfrule}{\mybox{by inference rule}}


% \renewcommand{\tin}{\mathsf{inn}}
% \renewcommand{\thead}{\cop{\mathsf{head}}}
% \renewcommand{\ttail}{\cop{\mathsf{tail}}}
% % start copattern color
% \newcommand{\shead}{\cop{\mathsf{head}}}
% \newcommand{\stail}{\cop{\mathsf{tail}}}
% \newcommand{\phead}{\cop{\mathsf{.head}}}
% \newcommand{\ptail}{\cop{\mathsf{.tail}}}
% \newcommand{\vecQ}{\cop{\vec Q}}
% \newcommand{\bDelta}{\boring{\Delta}}
% \newcommand{\bcovered}{\mathrel{\boring{\ccovered}}}
% \newcommand{\bA}{\boring{A}}
% \newcommand{\bC}{\boring{C}}
% \newcommand{\bder}{\mathrel{\boring{\vdash}}}
% \newcommand{\bcolon}{\mathrel{\boring{:}}}
% \newcommand{\blpar}{\boring{(}}
% \newcommand{\brpar}{\boring{)}}
% \renewcommand{\vbranches}{\boring{\mathit{branches}}}
% \renewcommand{\tout}{\cop{\mathsf{out}}}
% \renewcommand{\tforce}{\cop{\mathsf{force}}}
% \newcommand{\kw}[1]{{\bf\texttt{#1}}}
% \newcommand{\kwnewtype}{\kw{newtype}}
% \newcommand{\kwrecord}{\kw{record}}
% \newcommand{\kwwhere}{\kw{where}}
% \newcommand{\kwinstance}{\kw{instance}}
% \newcommand{\kwcoinductive}{\kw{coinductive}}
% \newcommand{\kwCoInductive}{\kw{CoInductive}}
% \newcommand{\kwCoFixpoint}{\kw{CoFixpoint}}
% \newcommand{\kwfield}{\kw{field}}
% \renewcommand{\tcase}{\mathbf{case}}
% \renewcommand{\tof}{\mathbf{of}}
% \renewcommand{\caseof}[3]{\tcase\,#1\,\tof~#2 \mathrel{\boldsymbol\To} #3}
% \newcommand{\tthead}{\cop{\texttt{head}}}
% \newcommand{\tttail}{\cop{\texttt{tail}}}
\newcommand{\ONE}{\mathbf{1}}

% DOES NOT WORK:
% rule with "boring" default text
\newcommand{\setboring}{\setbeamercolor{math text}{fg=grey}\setbeamercolor{math display}{fg=grey}}
\newcommand{\bru}[2]{\setboring\dfrac{\setboring#1}{\setboring#2}}

\renewcommand*\ttdefault{txtt} % for listing package

% \newcommand{\defHaskelllistings}{%
%   \lstset{%
%     language=Haskell,%
%     basicstyle=\ttfamily\small\color{darkdirtyblue},% \ttfamily
%     keywordstyle=\ttfamily\bfseries,% \underbar
%     identifierstyle=,%
%     commentstyle=\itshape,%
%     columns=flexible,%spaceflexible,% fixed,% flexible,%
%     showstringspaces=false,%
% %    xleftmargin=\codeindent,% defined below
%     breaklines=true,%
%     deletekeywords={succ,zero,head,tail,zipWith,Either,List},%
%     morekeywords={Set,Size,fun,cofun,pattern},% ,left,right,nil,cons
%     literate={\\}{{$\lambda$}}1 {->}{{$\rightarrow$~}}2
%              {<=}{{$\leq$~}}2 {<}{{$<$~}}1
% %     literate={map}{map~}4
% %       {even}{even~}5
% %       {odd}{odd~}4
%      }%
% }
%\defHaskelllistings

\title[Graded CBPV]{Graded Call-by-Push-Value}
\subtitle{Tracking (co)effects in a functional language}

\author[Abel]{
  Andreas Abel%\inst{1}
}
% \author{Andreas Abel\inst{1}
%   \and Brigitte Pientka\inst{2}
%   \and David Thibodeau\inst{2}
%   \and Anton Setzer\inst{3}
% }
%{F.~Author\inst{1} \and S.~Another\inst{2}}
% - Give the names in the same order as the appear in the paper.
% - Use the \inst{?} command only if the authors have different
%   affiliation.

\institute[] %Chalmers/GU/ENS Cachan] % (optional, but mostly needed)
{
  %\inst{1}%
  Department of Computer Science and Engineering\\
  Chalmers and University of Gothenburg, Sweden %\\[1ex]
}
%  \inst{1}%
%  Department of Computer Science\\
%  University of Somewhere
%  \and
%  \inst{2}%
%  Department of Theoretical Philosophy\\
%  University of Elsewhere}
%% - Use the \inst command only if there are several affiliations.
%% - Keep it simple, no one is interested in your street address.

\date[LFCS'25] % (optional, should be abbreviation of conference name)
{ LFCS seminar, U Edinburgh, Scotland \\
  17 June 2025
}
% - Either use conference name or its abbreviation.
% - Not really informative to the audience, more for people (including
%   yourself) who are reading the slides online

%\subject{Software Verification}
% This is only inserted into the PDF information catalog. Can be left
% out.



% If you have a file called "university-logo-filename.xxx", where xxx
% is a graphic format that can be processed by latex or pdflatex,
% resp., then you can add a logo as follows:

% \pgfdeclareimage[height=0.5cm]{university-logo}{university-logo-filename}
% \logo{\pgfuseimage{university-logo}}



% Delete this, if you do not want the table of contents to pop up at
% the beginning of each subsection:

%\AtBeginSubsection[]
%\AtBeginSection[]
%{
%  \begin{frame}<beamer>
%    \frametitle{Outline}
%    \tableofcontents[currentsection,currentsubsection]
%  \end{frame}
%}


% If you wish to uncover everything in a step-wise fashion, uncomment
% the following command:

%\beamerdefaultoverlayspecification{<+->}

\makeatletter
\def\MLine#1{\par\vspace{1ex}\hspace*{-\@totalleftmargin}\parbox{\textwidth}{\[#1\]}}
\makeatother
%\newenvironment{display}{\par\vspace{1ex}\hspace*{-\@totalleftmargin}\parbox{\textwidth}{\[}{\]}}


%\newcommand{\List}{\mathsf{List}}
\renewcommand{\Set}{\mathsf{Set}}
\newcommand{\Tree}{\mathsf{Tree}}
%\newcommand{\Prop}{\mathsf{Prop}}
%\newcommand{\Type}{\mathsf{Type}}
%\newcommand{\Size}{\mathsf{Size}}
%\newcommand{\tfix}{\mathsf{fix}}
%\newcommand{\Int}{\mathsf{Int}}
%\newcommand{\tnil}{\mathsf{nil}}
%\newcommand{\tcons}{\mathsf{cons}}
%\newcommand{\vas}{\mathit{as}}
%\newcommand{\ttuple}[1]{(#1)}

%\newcommand{\oford}{\of\tord} \newcommand{\oftype}{\of\ttype}
\newcommand{\oford}{}         \newcommand{\oftype}{}
\newenvironment{prg}{\begin{quotation}\begin{tabbing}}{\end{tabbing}\end{quotation}}

% COLORS
\newcommand{\cHead}{\color{darkblue}}
\newcommand{\cSub}{\color{brown}}
\newcommand{\cWhite}{\color{white}}
\newcommand{\cGray}{\color{gray}}
\newcommand{\cGreen}{\color{olivegreen}}
\newcommand{\cBrown}{\color{brown}}
\newcommand{\cBlack}{\color{black}}
\newcommand{\black}[1]{{\cBlack#1}}

\newcommand{\cAnn}{\color{red!80!black}}%purple darkblue
\newcommand{\cAside}{\color{gray}}
\newcommand{\cEnum}{\color{darkgreen}}
\newcommand{\cEm}{\cAnn} %\color{red}}
\newcommand{\cCo}{\cAnn} %\color{red}} % copattern color
\newcommand{\cop}[1]{{\cCo#1}}
\newcommand{\cApp}{\color{violet}}
\newcommand{\capp}[1]{{\cApp#1}}
\newcommand{\cFocus}{\color{darkgreen}}
\newcommand{\focus}[1]{{\cFocus#1}}
\newcommand{\cMath}{\usebeamercolor[fg]{math text}}
\newcommand{\cIdent}{\usebeamercolor[fg]{math text}}
\newcommand{\ident}[1]{{\cIdent#1}}
\newcommand{\cExp}{\cIdent}
\newcommand{\cBoring}{\color{grey}}
\newcommand{\boring}[1]{{\cBoring#1}}



\renewcommand{\rulename}[1]{#1}

% types
\newcommand{\cType}{\color{orange!60!black}}


\newcommand{\mlsays}[1]{}
% {\begin{frame}%[fragile=singleslide]
%   %\frametitle{Himself}
%   \begin{minipage}[c]{0.5\linewidth}
%    \includegraphics[height=0.9\textheight]{martin-loef-bw.png}
%   \end{minipage}% NO SPACE HERE!
%   \begin{minipage}[c]{0.5\linewidth}
%    \begin{center}
%      \bla \\[8ex]
%      \Huge #1
%    \end{center}
%   \end{minipage}
% \end{frame}
% }

\newcommand{\JOKE}[1]{} % No jokes.
\renewcommand{\emph}[1]{\textit{\cType#1}}


% \DeclareMathOperator*{\amp}{\&}
% \DeclareMathOperator*{\bigamp}{\scalerel*{\&}{\sum}}
% \usepackage{scalerel}

\begin{document}

\maketitle
%\begin{frame}
%  \titlepage
%\end{frame}

%\begin{frame}
%  \frametitle{Outline}
%  \tableofcontents
%  % You might wish to add the option [pausesections]
%\end{frame}

\section{Call-by-Push-Value}


% % Types
% \newcommand{\tTy}{\mathsf{Ty}}
% \newcommand{\Ty}[1]{\tTy^{#1}}
% \newcommand{\PTy}{\Ty+}
% \newcommand{\NTy}{\Ty-}
% \newcommand{\Cxt}{\mathsf{Cxt}}
% \newcommand{\RCxt}{\mathsf{RCxt}}
% \newcommand{\cempty}{\emptyset}
% \newcommand{\sempty}{\emptyset}
% \newcommand{\sext}[3]{#1,#3/#2}
% \newcommand{\extr}[2]{#1.#2}
% \newcommand{\ext}[3]{#1.#2{:}#3}
% \newcommand{\qext}[5]{#1#2.#3#4{:}#5}
% \newcommand{\gqext}[5]{\graybox{#1}#2.\graybox{#3}#4{:}#5}


\newcommand{\sumt}   [2]{\ensuremath{\Sigma_{#1} #2}}
\newcommand{\tupt}   [2]{\ensuremath{\mathop\otimes_{#1} #2}}
\newcommand{\rect}   [2]{\ensuremath{\Pi_{#1} #2}}
\newcommand{\sumty}  [3]{\ensuremath{\Sigma_{#1:#2} #3_{#1}}}
\newcommand{\tupty}  [3]{\ensuremath{\mathop\otimes_{#1:#2} #3_{#1}}}
\newcommand{\recty}  [3]{\ensuremath{\Pi_{#1:#2} #3_{#1}}}
\newcommand{\thunkty}[2]{\ensuremath{[#1]#2}}
\newcommand{\compty} [2]{\ensuremath{\langle#1\rangle#2}}


\begin{frame}%[fragile=singleslide]
  \frametitle{Call-by-Push-Value (CBPV)}
  %\vspace{-3ex}
  \begin{itemize}
  \item Value (positive) types $P$ %% ::= \forever N \mid \dotone \mid P \dottimes P' \mid \dotzero \mid P \dotplus P'$
\[
  \begin{array}{ll}
    o & \mbox{base type (e.g. integers)} \\
    \sumty i I P & \mbox{variants/sums with labels $i:I$} \\
    \tupty i I P & \mbox{tuples/vectors with indices $i:I$} \\
    \Box N       & \mbox{thunks}
  \end{array}
\]
  \item Computation (negative) types $N$ %% ::= \sometimes P \mid P \arrow N \mid N  $
\[
  \begin{array}{ll}
    \diamond P & \mbox{monadic type} \\
    P \To N & \mbox{function type (cbv)} \\
    \recty i I N & \mbox{records with fields $i:I$} \\
  \end{array}
\]
  \end{itemize}
\end{frame}

\renewcommand{\tcase}{\mathsf{case}}

\begin{frame}[fragile=singleslide]
  \frametitle{CBPV typing: values}
  %\vspace{-3ex}
  \begin{itemize}
  %\item A derivation of $P \in \Gamma$ is a (value) variable.
  \item A derivation of $\Gamma \vdash P$ is a value expression, $P \in \Gamma$ is a variable.
\begin{gather*}
  \nrux{\tin_i}{\Gamma \vdash P_i}{\Gamma \vdash \sumt I P}{i{:}I}
 \qquad
  \nru{\mathsf{tup}}{\forall i{:}I,\, \Gamma \vdash P_i}{\Gamma \vdash \tupt I P}
\\
  \nru{\tvar}{P \in \Gamma}{\Gamma \vdash P}
\qquad
  \nru{\tthunk}{\Gamma \vdash N}{\Gamma \vdash \Box N}
\end{gather*}
  \item A derivation of $\Gamma \vdash N$ is a computation expression.

  \item Computing with values:
\begin{gather*}
  \nru{\tcase}{\Gamma \vdash \sumt I P \quad \forall i{:}I,\, \Gamma.P_i \vdash N}{\Gamma \vdash N}
\qquad
  \nru{\tsplit}{\Gamma \vdash \tupt I P \quad \Gamma.\overline{P_i}^{i:I} \vdash N}{\Gamma \vdash N}
  \\
  \nru{\tlet}{\Gamma \vdash P \quad \Gamma.P \vdash N}{\Gamma \vdash N}
\qquad
  \nru{\tforce}{\Gamma \vdash \Box N}{\Gamma \vdash N}
\end{gather*}
  \end{itemize}
\end{frame}

\newcommand{\trecord}{\mathsf{record}}
\newcommand{\tproj}{\mathsf{proj}}

\begin{frame}[fragile=singleslide]
  \frametitle{CBPV typing: computations}
\begin{itemize}
  \item Computation introductions:
\begin{gather*}
  \nru{\treturn}{\Gamma \vdash P}{\Gamma \vdash \diamond P}
\qquad
  \nru{\lambda}{\Gamma.P \vdash N}{\Gamma \vdash P \To N}
\qquad
  \nru{\trecord}{\forall i{:}I,\, \Gamma \vdash N_i}{\Gamma \vdash \rect I N}
\end{gather*}
  \item Computation eliminations:
\begin{gather*}
  \nru{\tbind}{\Gamma \vdash \diamond P \quad \Gamma.P \vdash N}{\Gamma \vdash N}
\qquad
  \nru{\tapp}{\Gamma \vdash P \To N \quad \Gamma \vdash P}{\Gamma \vdash N}
\\
  \nru{\tproj_i}{\Gamma \vdash \rect I N}{\Gamma \vdash N_i}
\end{gather*}
  \end{itemize}
\end{frame}

\newcommand{\tprint}{\mathsf{print}}

\begin{frame}%[fragile=singleslide]
  \frametitle{CBPV: effect}
  %\vspace{-3ex}
  \begin{itemize}
  \item Let $\top = \rect \emptyset \_$ be the unit type.
  \item Simple effect: print string literals.
  \[
     \nru{\tprint\,s}{}{\Gamma \vdash \diamond \top}
  \]
  \item Example (using $\tbind$ infix):
  \[f = \tprint\,\texttt{"function"}\,\tbind\,\_. \lambda x.\tprint\,\texttt{"argument"}\,\tbind\,\_.\tret\,(\tvar\,x)\]
  \item CBPV effects are lazy, can only be observed at value types.
  \item $f$ by itself prints nothing; $\tapp\,f\,\_$ prints \texttt{"function"} and \texttt{"argument"}.
  \end{itemize}
\end{frame}

\begin{frame}%[fragile=singleslide]
  \frametitle{Monad and monad algebras}
\renewcommand{\T}{\diamond}
\[
\xymatrix@C=12ex{
  \T\,A     \ar[r]^{\treturn} \ar@{=}[dr]^{\tid}
& \T\,(\T\,A) \ar[d]^{\tjoin}
& \T\,(\T\,(\T\,A)) \ar[l]_{\tjoin} \ar[d]^{\tmap\,\tjoin}
\\
& \T\,A       \ar@{=}[dr]^{\tid}
& \T\,(\T\,A) \ar[l]_{\tjoin}
\\
&
& \T\,A \ar[u]_{\tmap\,\treturn}
}
\]
\[
\xymatrix@C=10ex{
  B     \ar[r]^{\treturn} \ar@{=}[dr]^{\tid}
& \T\,B \ar[d]^{\trun}
& \T\,(\T\,B) \ar[l]_{\tjoin} \ar[d]^{\tmap\,\trun}
\\
& B
& \T\,B \ar[l]_{\trun}
}
\]
  % \vspace{-3ex}
  % \begin{itemize}
  % \item
  % \end{itemize}
\end{frame}

\newcommand{\tduplicate}{\mathsf{duplicate}}
\newcommand{\tdisplay}{\mathsf{display}}
\newcommand{\texpose}{\mathsf{expose}}

\begin{frame}%[fragile=singleslide]
  \frametitle{CBPV: algebraic semantics}
  %\vspace{-3ex}
  \begin{itemize}
  %\item Interpret both $P$ and $N$ as objects in a cartesian closed category with coproducts.
  \item Assume a cartesian closed category with coproducts.
  \item Contexts $P_1\ldots P_n$ are products $P_1 \times \ldots \times P_n$ and $\Gamma \vdash A$ homsets.
  \item Interpret $\diamond$ as a strong monad; $\Box$ as identity (for now).
  %, $\Box$ as a strong comonad.
  \item Interpret each $N$ as monad algebra $\trun_N : \diamond N \to N$.
  \[
    \begin{array}{lcl}
      \trun_{\diamond P} & : & \diamond (\diamond P) \to \diamond P \\
      \trun_{\diamond P} & = & \tjoin \\[1ex]
      \trun_{P \To N} & : & \diamond (P \To N) \to P \To N \\
      \trun_{P \To N}\, m & = & \lambda p.\, \trun_N\,(\tmap\, (\lambda f.\, \tapp\,f\,p)\, m) \\[1ex]
      \trun_{\rect I N} & : & \diamond (\rect I N) \to \rect I N \\
      \trun_{\rect I N}\,m & = & \trecord\,i.\,\trun_N\,(\tmap\, \tproj_i\, m)
    \end{array}
  \]
  \item This interprets binding in $N$ (going via $\diamond N$):
\[
  \nru{\tbind}{\Gamma \vdash \diamond P \quad \Gamma.P \vdash N}{\Gamma \vdash N}
\]
  \end{itemize}
\end{frame}

\section{Effect-graded CBPV}

\newcommand{\tsub}{\mathsf{sub}}

\begin{frame}%[fragile=singleslide]
  \frametitle{Graded effects}
  %\vspace{-3ex}
  \begin{itemize}
  \item $\Gamma \vdash N \mid e$ where $e$ is drawn from an effect pomonoid.
  \item Example: keep track of the length of the output.
  \[
     \nrux{\tprint\,s}{}{\Gamma \vdash \diamond \top \mid n}{n = |s|}
  \]
  \item Effects accumulate.
\[
  \nru{\treturn}{\Gamma \vdash P}{\Gamma \vdash \diamond P \mid 0}
\qquad
  \nru{\tbind}{\Gamma \vdash \diamond P \mid e_1 \quad \Gamma.P \vdash N \mid e_2}{\Gamma \vdash N \mid e_1 + e_2}
\]
  \item Effects get frozen and subsumed.
\[
  \nru{\tthunk}{\Gamma \vdash N \mid e}{\Gamma \vdash [e]N}
\qquad
  \nru{\tforce}{\Gamma \vdash [e]N}{\Gamma \vdash N \mid e}
\qquad
  \nrux{\tsub}{\Gamma \vdash N \mid e}{\Gamma \vdash N \mid e'}{e \leq e'}
\]
  \end{itemize}
\end{frame}

\newcommand{\bu}{\bullet}
\newcommand{\Char}{\mathsf{Char}}

\begin{frame}[fragile=singleslide]
  \frametitle{Graded monads}
  \vspace{-3ex}
\renewcommand{\T}[1][]{\diamond_{#1}}
\[
  \begin{array}{lcl@{\qquad}rcl}
    \multicolumn 3 r {} %{\mbox{Example:}}
      & \T[n] A & = & \Char^n \times A \\
    \treturn & : & A \to \T[\Ge] A
      & a & \mapsto & (\mathtt{""},a) \\
    \tjoin  & : & \T[e_1](\T[e_2] A) \to \T[e_1 \bu e_2] A
      & (s_1,(s_2,a)) & \mapsto & (s_1 + s_2, a) \\
    \T[e \leq e'] & : & \T[e] A \to \T[e'] A \\
  \end{array}
\]
\[
\xymatrix@C=8ex{
  \T[e]\,A     \ar[r]^{\treturn} \ar@{=}[dr]^{\tid}
& \T[\Ge]\,(\T[e]\,A) \ar[d]^{\tjoin}
\\
& \T[e]\,A
}
\quad
\xymatrix@C=10ex{
  \T[e]\,A     \ar[r]^{\tmap\,\treturn} \ar@{=}[dr]^{\tid}
& \T[e]\,(\T[\Ge]\,A) \ar[d]^{\tjoin}
\\
& \T[e]\,A
}
\]
\[
\xymatrix@C=8ex{
  \T[e_1 \bu e_2]\,(\T[e_3]\,A) \ar[d]^{\tjoin}
& \T[e_1]\,(\T[e_2]\,(\T[e_3]\,A)) \ar[l]_{\tjoin} \ar[d]^{\tmap\,\tjoin}
\\
  \T[e_1 \bu e_2 \bu e_3]\,A
& \T[e_1]\,(\T[e_2 \bu e_3]\,A) \ar[l]_{\tjoin}
}
\]
\end{frame}


\newcommand{\run}[1][]{\mathsf{run}_{#1}}

\begin{frame}[fragile=singleslide]
  \frametitle{Graded monad algebras}
  \vspace{-3ex}
\renewcommand{\T}[1][]{\diamond_{#1}}
\[
  \begin{array}{lcl}
    B_{e \leq e'} & : & B_e \to B_{e'} \\
    \run[B] & : & \T[e_1]\,B_{e_2} \to B_{e_1 \bu e_2} \\
  \end{array}
\]
\[
\xymatrix@C=10ex{
  B_e     \ar[r]^{\treturn} \ar@{=}[dr]^{\tid}
& \T[\Ge]\,B_e \ar[d]^{\run[]}
\\
& B_e
}
\qquad
\xymatrix@C=10ex{
  \T[e_1 \bu e_2]\,B_{e_3} \ar[d]^{\run[]}
& \T[e_1]\,(\T[e_2]\,B_{e_3}) \ar[l]_{\tjoin} \ar[d]^{\tmap\,\run[]}
\\
  B_{e_1 \bu e_2 \bu e_3}
& \T[e_1]\,B_{e_2 \bu e_3} \ar[l]_{\run[]}
}
\]
\[
  \begin{array}{lcl}
    (\diamond P)_e & = & \diamond_e P \\
    (P \To N)_e    & = & P \To N_e \\
    (\tupty i I N)_e  & = & \tupt {i:I} {(N_i)_e} \\
  \end{array}
\]
\end{frame}



\begin{frame}%[fragile=singleslide]
  \frametitle{Graded effect semantics}
  % \vspace{-3ex}
  \begin{itemize}
  \item Negative types $N$ are interpreted as functors from effects to objects.
  \item Positive types $P$ are interpreted as objects; $[e]N = N_e$.
  \item $\Gamma \vdash P$ is $\Gamma \to P$.
  \item $\Gamma \vdash N \mid e$ is $\Gamma \to N_e$.
\[
  \nru{\treturn}{\Gamma \to P}{\Gamma \to \diamond_\Ge P}
\qquad
  \nru{\tbind}{\Gamma \to \diamond_{e_1} P \quad \Gamma \times P \to N_{e_2}}{\Gamma \to N_{e_1 \bu e_2}}
\]
  $\tbind$ justified by $\Gamma \to \Gamma \times \diamond_{e_1} P \to \diamond_{e_1}(\Gamma \times P) \to \diamond_{e_1}N_{e_2} \to N_{e_1 \bu e_2}$.
  \end{itemize}
\end{frame}

\section{Coeffect CBPV}

\newcommand{\textract}{\mathsf{extract}}
%\newcommand{\tdisplay}{display}

\begin{frame}%[fragile=singleslide]
  \frametitle{Comonads and comonad coalgebras}
\renewcommand{\D}{\Box}
\[
\xymatrix@C=10ex{
\D\,B  \ar@{=}[dr]
  & \D\,(\D\,B) \ar[l]_{\textract} \ar[r]^{\tdisplay}
  & \D\,(\D\,(\D\,B))
\\
  & \D\,B \ar[u]_{\tdisplay} \ar[r]^{\tdisplay} \ar@{=}[dr]
  & \D\,(\D\,B) \ar[u]_{\tmap\,\tdisplay} \ar[d]^{\tmap\,\textract}
\\
  &
  & \D\,B
}
\]
\[
\xymatrix@C=10ex{
A  \ar@{=}[dr]
  & \D\,A \ar[l]_{\textract} \ar[r]^{\tdisplay}
  & \D\,(\D\,A)
\\
  & A \ar[u]_{\texpose} \ar[r]^{\texpose}
  & \D\,A \ar[u]_{\tmap\,\texpose}
}
\]
\end{frame}


\begin{frame}[fragile=singleslide]
  \frametitle{Comonadic semantics}
  %\vspace{-3ex}
  \begin{itemize}
  \item Judgements $P_1 \ldots P_n \vdash A$ are homsets $P_1 \times \ldots \times P_n \to A$.
  \item $\Box$ is interpreted as monoidal comonad ($\diamond$ as identity).
  \item Positive types $P$ are $\Box$-coalgebras (negative types $N$ just objects).
\[
\begin{array}{lcl}
  \texpose_{\Box N} & : & \Box N \to \Box (\Box N) \\
  \texpose_{\Box N} & = & \tdisplay \\[1ex]
  \texpose_{\sumt I P} & : & \sumt I P \to \Box (\sumt I P) \\
  \texpose_{\sumt I P} (i, v) & = & \tmap\,\tin_i\, (\texpose_{P_i} v) \\[1ex]
  % \texpose_{\sumt I P} & = & \left [ \tmap\,\tin_i \comp \texpose_{P_i} \right ]_{i:I} \\[1ex]
  \texpose_{\tupt I P} & : & \tupt I P \to \Box(\tupt I P) \\
  \texpose_{\tupt I P} (v_i)_{i:I} & = & \tzip\, (\texpose_{P_i} v_i)_{i:I} \\[1ex]
  \tzip & : & \tupt {i:I} (\Box {B_i}) \to \Box (\tupt I B) \\[1ex]
  \texpose_\Gamma & : & \Gamma \to \Box \Gamma \\[1ex]
\end{array}
\]
  \end{itemize}
\end{frame}


\begin{frame}%[fragile=singleslide]
  \frametitle{}
  % \vspace{-3ex}
  \begin{itemize}
  \item Thunking needs monoidality.
\[
  \nru{\tthunk}{\Gamma \to N}{\Gamma \to \Box N}
\qquad
  \nru{\tforce}{\Gamma \to \Box N}{\Gamma \to N}
\]
\[
  \begin{array}{lcl}
    \tthunk & : & \Gamma \to \Box \Gamma \to \Box N \\
    \tforce & : & \Gamma \to \Box N \to N \\
  \end{array}
\]
  \end{itemize}
\end{frame}


\section{Coeffect-graded CBPV}

\newcommand{\lolli}{\multimap}



\begin{frame}[fragile=singleslide]
  \frametitle{Graded comonad}
\vspace{-2ex}
\renewcommand{\D}[1][]{\Box_{#1}}
\[
  \begin{array}{lcl}
    \textract & : & \D[1] B \to B \\
    \tdisplay & : & \D[qr] B \to \D[q](\D[r]B) \\
  \end{array}
\]
\[
\xymatrix@C=10ex{
\D[r]\,B  \ar[d]_{\tdisplay} \ar[r]^{\tdisplay} \ar@{=}[dr]
  & \D[r]\,(\D[1]\,B)  \ar[d]^{\tmap\,\textract}
\\
\D[1]\,(\D[r]\,B) \ar[r]^{\textract}
  & \D[r]\,B
}
\]
\[
\xymatrix@C=12ex{
\D[qrs]\,B \ar[d]_{\tdisplay} \ar[r]^{\tdisplay}
  & \D[qr]\,(\D[s]\,B) \ar[d]^{\tdisplay}
\\
\D[q]\,(\D[rs]\,B) \ar[r]^{\tmap\,\tdisplay}
  & \D[q]\,(\D[r]\,(\D[s]\,B))
}
\]
\end{frame}


\begin{frame}[fragile=singleslide]
  \frametitle{Graded comonad coalgebras}
  \vspace{-3ex}
\renewcommand{\D}[1][]{\Box_{#1}}
\[
  \begin{array}{lcl}
    A_{q \leq r} & : & A_q \to A_r \\
    \texpose & : & A_{qr} \to \D[q]A_r \\
  \end{array}
\]
\[
\xymatrix@C=10ex{
A_r  \ar[d]_{\texpose} \ar@{=}[dr]
\\
\D[1]\,A_r \ar[r]^{\textract}
  & A_r
}
\qquad
\xymatrix@C=12ex{
A_{qrs} \ar[d]_{\texpose} \ar[r]^{\texpose}
  & \D[qr]\,A_s \ar[d]^{\tdisplay}
\\
\D[q]\,A_{rs} \ar[r]^{\tmap\,\texpose}
  & \D[q]\,(\D[r]\,A_s)
}
\]
\end{frame}

\newcommand{\tdrop}{\mathsf{drop}}
\newcommand{\I}{\mathsf{I}}

\begin{frame}[fragile=singleslide]
  \frametitle{Semiring-graded comonads}
  %\vspace{-3ex}
  \begin{itemize}
  \item Switch to a closed symmetric monoidal category.
  \item Additive commutative monoid on grades maps to monoidal structure.
\[
  \begin{array}{lcl}
    \tdrop & : & \Box_0 B \to \I \\
    \tduplicate & : & \Box_{r+s} B \to \Box_r B \otimes \Box_s B \\
  \end{array}
\]
  \item Extends to coalgebras.
\[
  \begin{array}{lcl}
    \tdrop_A & : & A_0 \to \I \\
    \tduplicate_A & : & A_{r+s} \to A_r \otimes A_s \\
  \end{array}
\]
  \end{itemize}
\end{frame}

\begin{frame}%[fragile=singleslide]
  \frametitle{Semantics of positive types}
  %\vspace{-3ex}
  %\begin{itemize}
  %\item
Positive types $P$ are interpreted as graded $\Box$-coalgebras.
\[
  \begin{array}{lcl}
    (\Box N)_r & = & \Box_r N \\
    (\sumt I P)_r & = & \sumt {i:I} {(P_i)_r} \\
    (\tupt I P)_r & = & \tupt {i:I} {(P_i)_r} \\
  \end{array}
\]
\[
  \begin{array}{lcl}
    \texpose_{\Box N} & : & (\Box N)_{qr} \to \Box_q (\Box N)_r \\
    \texpose & = & \tdisplay \\[1ex]
    \texpose_{\sumt I P} & : & (\sumt I P)_{qr} \to \Box_q (\sumt I P)_r\\
    \texpose_{\sumt I P} (i, v) & = & \tmap\, \tin_i\, (\texpose_{P_i} v) \\[1ex]
    \texpose_{\tupt I P} & : & (\tupt I P)_{qr} \to \Box_q (\tupt I P)_r \\
    \texpose_{\tupt I P} (v_i)_{i:I} & = & \tzip\,(\texpose_{P_i} v_i)_{i:I} \\
  \end{array}
\]
%  \end{itemize}
\end{frame}


\begin{frame}[fragile=singleslide]
  \frametitle{Graded contexts}
  %\vspace{-3ex}
  \begin{itemize}
  \item Contexts $\Gamma$ come with grading contexts $\gamma$ of the same length.
  \item $r_1P_1\dots r_nP_n = (r_1\dots r_n)(P_1 \dots P_n)$.
  \item Grading contexts $\gamma$ (``vectors'') form a \emph{left module} to grades.
  \item Write $\mathsf{e}_x$ for the unit vector $e_x(y) = 1$ iff $x=y$ else $0$.
  \item Contexts $\Gamma$ are interpreted as \emph{structured} graded $\Box$-coalgebras $\Gamma_\gamma$.
\[
  \begin{array}{lcl}
    \texpose_\Gamma& : & \Gamma_{r\gamma} \to \Box_r \Gamma_\gamma \\
    \tdrop_\Gamma & : & \Gamma_0 \to \I \\
    \tduplicate_\Gamma & : & \Gamma_{\gamma+\delta} \to \Gamma_\gamma \otimes \Gamma_\delta \\
    \tduplicate_\Gamma & : & \Gamma_{\sum_I \gamma} \to \tupt {i:I} \Gamma_{\gamma_i} \\
  \end{array}
\]
  \end{itemize}
\end{frame}


% \begin{frame}%[fragile=singleslide]
%   \frametitle{Coeffect-graded CBPV}
%   %\vspace{-3ex}
%   \begin{itemize}
%   \item Negative types $N ::= \langle r \rangle P \mid rP \lolli N \mid \recty i I N$.
%   \end{itemize}
% \end{frame}

% \begin{frame}[fragile=singleslide]
%   \frametitle{Coeffect-graded CBPV}
%   %\vspace{-3ex}
%   \begin{itemize}
%   %\item A derivation of $P \in \Gamma$ is a (value) variable.
%   \item Negative types $N ::= \langle r \rangle P \mid rP \lolli N \mid \recty i I N$.
%   \item $\gamma\Gamma \vdash P$ and $\gamma\Gamma \vdash N$.
% \begin{gather*}
%   \nrux{\tin_i}{\gamma\Gamma \vdash P_i}{\gamma\Gamma \vdash \sumt I P}{i{:}I}
%  \qquad
%   \nru{\mathsf{tup}}{\forall i{:}I,\, \gamma_i\Gamma \vdash P_i}{(\sum_I \gamma) \Gamma \vdash \tupt I P}
% \\
%   \nru{\tvar}{x : P \in \Gamma}{\mathsf{e}_x\Gamma \vdash P}
% \qquad
%   \nru{\tthunk}{\gamma\Gamma \vdash N}{\gamma\Gamma \vdash \Box N}
% \end{gather*}
%   \item Computing with values:
% \begin{gather*}
%   \nru{\tcase}{\delta\Gamma \vdash \sumt I P \quad \forall i{:}I,\, \gamma\Gamma.rP_i \vdash N}{(\gamma + r\delta)\Gamma \vdash N}
% \qquad
%   \nru{\tsplit}{\delta\Gamma \vdash \tupt I P \quad \gamma\Gamma.\overline{rP_i}^{i:I} \vdash N}{(\gamma + r\delta)\Gamma \vdash N}
%   \\
%   \nru{\tlet}{\delta\Gamma \vdash P \quad \gamma\Gamma.rP \vdash N}{(\gamma + r\delta)\Gamma \vdash N}
% \qquad
%   \nru{\tforce}{\gamma\Gamma \vdash \Box N}{\gamma\Gamma \vdash N}
% \end{gather*}
%   \end{itemize}
% \end{frame}


% \begin{frame}[fragile=singleslide]
%   \frametitle{CBPV typing: computations}
% \begin{itemize}
%   \item Computation introductions:
% \begin{gather*}
%   \nru{\treturn}{\gamma\Gamma \vdash P}{r\gamma\Gamma \vdash \langle r \rangle P}
% \qquad
%   \nru{\lambda}{\gamma\Gamma.rP \vdash N}{\gamma\Gamma \vdash rP \lolli N}
% \qquad
%   \nru{\trecord}{\forall i{:}I,\, \gamma\Gamma \vdash N_i}{\gamma\Gamma \vdash \rect I N}
% \end{gather*}
%   \item Computation eliminations:
% \begin{gather*}
%   \nru{\tbind}{\delta \Gamma \vdash \langle r \rangle P \quad \gamma\Gamma.r P \vdash N}{(\gamma + \delta)\Gamma \vdash N}
% \qquad
%   \nru{\tapp}{\gamma\Gamma \vdash rP \To N \quad \delta\Gamma \vdash P}{(\gamma + r\delta)\Gamma \vdash N}
% \\
%   \nru{\tproj_i}{\gamma\Gamma \vdash \rect I N}{\gamma\Gamma \vdash N_i}
% \end{gather*}
%   \end{itemize}
% \end{frame}

\begin{frame}[fragile=singleslide]
  \frametitle{Coeffect-graded CBPV}
  %\vspace{-3ex}
  \begin{itemize}
  %\item A derivation of $P \in \Gamma$ is a (value) variable.
  \item Negative types $N ::= \langle r \rangle P \mid rP \lolli N \mid \recty i I N$.
  \item $\gamma\Gamma \vdash P$ and $\gamma\Gamma \vdash N$.
\begin{gather*}
  \nru{\tvar}{x : P \in \Gamma}{\mathsf{e}_x\Gamma \vdash P}
\qquad
  \nru{\tlet}{\delta\Gamma \vdash P \quad \gamma\Gamma.rP \vdash N}{(\gamma + r\delta)\Gamma \vdash N}
\\
  \nrux{\tin_i}{\gamma\Gamma \vdash P_i}{\gamma\Gamma \vdash \sumt I P}{i{:}I}
 \qquad
  \nru{\tcase}{\delta\Gamma \vdash \sumt I P \quad \forall i{:}I,\, \gamma\Gamma.rP_i \vdash N}{(\gamma + r\delta)\Gamma \vdash N}
\\
  \nru{\mathsf{tup}}{\forall i{:}I,\, \gamma_i\Gamma \vdash P_i}{(\sum_I \gamma) \Gamma \vdash \tupt I P}
\qquad
  \nru{\tsplit}{\delta\Gamma \vdash \tupt I P \quad \gamma\Gamma.\overline{rP_i}^{i:I} \vdash N}{(\gamma + r\delta)\Gamma \vdash N}
\\
  \nru{\tthunk}{\gamma\Gamma \vdash N}{\gamma\Gamma \vdash \Box N}
\qquad
  \nru{\tforce}{\gamma\Gamma \vdash \Box N}{\gamma\Gamma \vdash N}
\end{gather*}
  \end{itemize}
\end{frame}


\begin{frame}[fragile=singleslide]
  \frametitle{CBPV typing: computations}
%\begin{itemize}
%  \item Computation introductions:
\begin{gather*}
  \nru{\treturn}{\gamma\Gamma \vdash P}{r\gamma\Gamma \vdash \langle r \rangle P}
\qquad
  \nru{\tbind}{\delta \Gamma \vdash \langle r \rangle P \quad \gamma\Gamma.r P \vdash N}{(\gamma + \delta)\Gamma \vdash N}
\\[1ex]
  \nru{\lambda}{\gamma\Gamma.rP \vdash N}{\gamma\Gamma \vdash rP \lolli N}
\qquad
  \nru{\tapp}{\gamma\Gamma \vdash rP \lolli N \quad \delta\Gamma \vdash P}{(\gamma + r\delta)\Gamma \vdash N}
\\[1ex]
  \nru{\trecord}{\forall i{:}I,\, \gamma\Gamma \vdash N_i}{\gamma\Gamma \vdash \rect I N}
\qquad
  \nru{\tproj_i}{\gamma\Gamma \vdash \rect I N}{\gamma\Gamma \vdash N_i}
\end{gather*}
%  \end{itemize}
\end{frame}

\begin{frame}%[fragile=singleslide]
  \frametitle{Semantics of coeffect-grading}
  %\vspace{-3ex}
  \begin{itemize}
  \item Computations $\gamma\Gamma \vdash N$ are morphisms $\Gamma_\gamma \to N$.
  \item Values $\gamma\Gamma \vdash P$ are morphisms families $\Gamma_{r\gamma} \to P_r$ natural in $r$.
\begin{gather*}
  \nru{\tvar}{x : P \in \Gamma}{\mathsf{e}_x\Gamma \vdash P}
\qquad
  \nru{\tlet}{\delta\Gamma \vdash P \quad \gamma\Gamma.rP \vdash N}{(\gamma + r\delta)\Gamma \vdash N}
\end{gather*}
\[
  \begin{array}{lcl}
    \tvar & : & (\Gamma.P.\Delta)_{s(\vec 0.1.\vec 0)}
       \to \Gamma_{\vec 0} \otimes P_s \otimes \Delta_{\vec 0}
       \to P_s \\
    \tlet & : & \Gamma_{\gamma + r\delta}
       \to \Gamma_{\gamma} \otimes \Gamma_{r\delta}
       \to \Gamma_{\gamma} \otimes P_{r}
       \to (\Gamma.P)_{\gamma.r}
       \to N \\
  \end{array}
\]
\[
  \nru{\tthunk}{\gamma\Gamma \vdash N}{\gamma\Gamma \vdash \Box N}
\qquad
  \nru{\tforce}{\gamma\Gamma \vdash \Box N}{\gamma\Gamma \vdash N}
\]
\[
  \begin{array}{lcl}
    \tthunk & : & \Gamma_{r\gamma} \to \Box_r\Gamma_\gamma \to \Box_rN = (\Box N)_r \\
    \tforce & : & \Gamma_\gamma = \Gamma_{1\gamma} \to (\Box N)_1 = \Box_1 N \to N \\
  \end{array}
\]
  \end{itemize}
\end{frame}


\begin{frame}%[fragile=singleslide]
  \frametitle{Semantics ctd.}
  %\vspace{-3ex}
  \begin{itemize}
  \item Semantics of monadic computations.
\[
  \nru{\treturn}{\gamma\Gamma \vdash P}{r\gamma\Gamma \vdash \langle r \rangle P}
\qquad
  \nru{\tbind}{\delta \Gamma \vdash \langle r \rangle P \quad \gamma\Gamma.r P \vdash N}{(\gamma + \delta)\Gamma \vdash N}
\]
\[
  \begin{array}{lcl}
    \treturn & : & \Gamma_{r\gamma} \to P_r \to \diamond P_r = \langle r \rangle P \\
    \tbind & : & \Gamma_{\gamma + \delta}
      \to \Gamma_\gamma \otimes \Gamma_\delta
      \to \Gamma_\gamma \otimes (\diamond P_r)
      \to \diamond (\Gamma_\gamma \otimes P_r)
      \to \diamond N
      \to N
  \end{array}
\]
\item Semantics of functions.
\[
  \nru{\lambda}{\gamma\Gamma.rP \vdash N}{\gamma\Gamma \vdash rP \lolli N}
\qquad
  \nru{\tapp}{\gamma\Gamma \vdash rP \lolli N \quad \delta\Gamma \vdash P}{(\gamma + r\delta)\Gamma \vdash N}
\]
\[
  \begin{array}{lcl}
    \tapp & : & \Gamma_{\gamma + r\delta}
       \to \Gamma_\gamma \otimes \Gamma_{r\delta}
       \to \Gamma_\gamma \otimes P_r
       \to N
  \end{array}
\]
  \end{itemize}
\end{frame}

\section{Fully graded CBPV}

\begin{frame}%[fragile=singleslide]
  \frametitle{Fully graded CBPV}
  \begin{itemize}
  \item Judgements $\gamma\Gamma \vdash P$ and $\gamma\Gamma \vdash N \mid e$.
  \item Interpreted as $\Gamma_{r\gamma} \to P_r$ and $\Gamma_\gamma \to N_e$.
  \item Monad $(\langle r \rangle P)_e = \diamond_e P_r$ and comonad $([e]N)_r = \Box_r N_e$.
  \item No commutative laws needed!
  \item No interaction between effects and coeffects in the framework.
  \end{itemize}
\end{frame}



% \begin{frame}[c]%[fragile=singleslide]
%   \frametitle{Formalization}
%   \begin{center}
% Agda code \\[4ex]

% \url{https://github.com/andreasabel/ipl/}
%   \end{center}
% \end{frame}



\begin{frame}%[fragile=singleslide]
\frametitle{References}
\begin{itemize}
\begin{small}

\item
\textbf{Paul Blain Levy}:
Call-by-push-value: Decomposing call-by-value and call-by-name.
HOSC 19(4), 2006.

\end{small}
\end{itemize}
\end{frame}


% \begin{frame}%[fragile=singleslide]
%   \frametitle{Future work}
%   %\vspace{-3ex}
%   \begin{itemize}
%   \item
%   \end{itemize}
% \end{frame}


%%%%%%%%%%%%%%%%%%%%%%%%%%%%%%%%%%%%%%%%%%%%%%%%%%%%%%%%%%%%%%%%%%%%%%
%%%%%%%%%%%%%%%%%%%%%%%%%%%%%% END DOC %%%%%%%%%%%%%%%%%%%%%%%%%%%%%%%
%%%%%%%%%%%%%%%%%%%%%%%%%%%%%%%%%%%%%%%%%%%%%%%%%%%%%%%%%%%%%%%%%%%%%%

% \bibliography{short}

\end{document}


\begin{frame}%[fragile=singleslide]
  \frametitle{}
  \vspace{-3ex}
  \begin{itemize}
  \item
  \end{itemize}
\end{frame}


\begin{frame}%[fragile=singleslide]
  \frametitle{}
  \vspace{-3ex}
  \begin{itemize}
  \item
  \end{itemize}
\end{frame}


\begin{frame}%[fragile=singleslide]
  \frametitle{}
  \vspace{-3ex}
  \begin{itemize}
  \item
  \end{itemize}
\end{frame}


\begin{frame}%[fragile=singleslide]
  \frametitle{}
  \vspace{-3ex}
  \begin{itemize}
  \item
  \end{itemize}
\end{frame}


\begin{frame}%[fragile=singleslide]
  \frametitle{}
  \vspace{-3ex}
  \begin{itemize}
  \item
  \end{itemize}
\end{frame}


\begin{frame}%[fragile=singleslide]
  \frametitle{}
  \vspace{-3ex}
  \begin{itemize}
  \item
  \end{itemize}
\end{frame}


\begin{frame}%[fragile=singleslide]
  \frametitle{}
  \vspace{-3ex}
  \begin{itemize}
  \item
  \end{itemize}
\end{frame}


\begin{frame}%[fragile=singleslide]
  \frametitle{}
  \vspace{-3ex}
  \begin{itemize}
  \item
  \end{itemize}
\end{frame}


\begin{frame}%[fragile=singleslide]
  \frametitle{}
  \vspace{-3ex}
  \begin{itemize}
  \item
  \end{itemize}
\end{frame}

\end{document}




%%% Local Variables:
%%% mode: latex
%%% TeX-master: t
%%% End:
